% Part: sets-functions-relations
% Chapter: functions
% Section: inverses

\documentclass[../../../include/open-logic-section]{subfiles}

\begin{document}

\olfileid{sfr}{fun}{inv}
\olsection{函数的逆}

\begin{explain}
我们通常把函数理解为映射。于是，关于函数，一个显而易见的问题是：映射能否被“反转”。例如，后继函数 $f(x) = x + 1$ 就是可以反转的，因为函数 $g(y) = y - 1$ “撤销”了 $f$ 的作用。

但我们必须小心。虽然 $g$ 的定义在 $\Int \to \Int$ 上确实定义了一个函数，但它并没有在 $\Nat \to \Nat$ 上定义一个\emph{函数}，因为 $g(0) \notin \Nat$。所以，即使在简单的情形中，一个函数能否被反转也并非完全显然；它可能依赖于定义域和陪域的选择。

函数的逆这一概念使上述想法得到了精确表述。
\end{explain}

\begin{defn}
给定函数 $f \colon A \to B$ 与 $g \colon B \to A$，如果对所有 $x \in A$ 和 $y \in B$ 都有 $f(g(y)) = y$ 且 $g(f(x)) = x$，则称 $g$ 是 $f$ 的一个\emph{逆}。
\end{defn}

如果 $f$ 有一个逆~$g$，我们通常将其记作 $f^{-1}$ 而非~$g$。

\begin{explain}
现在我们来探讨函数何时有逆。对于函数 $f\colon A \to B$，一个看似不错的逆的候选是 $g\colon B \to A$，它被“定义”为
\[
g(y) = \text{“那个” $x$ 使得 $f(x) = y$。}
\]
但“定义”（以及“那个”）周围的引号暗示着这并非一个真正的定义。至少，它并非在完全一般的意义下总能行得通。因为，要让这个定义确定一个函数，对于每个 $y \in B$，必须恰好存在一个 $x$ 使得 $f(x) = y$——亦即 $g$ 的输出必须是唯一确定的。此外，它对每个 $y \in B$ 都必须有定义。如果存在 $x_1, x_2 \in A$ 满足 $x_1 \neq x_2$ 但 $f(x_1) = f(x_2)$，那么 $g(y)$ 对于 $y = f(x_1) = f(x_2)$ 就不是唯一确定的。而如果根本不存在 $x$ 使得 $f(x) = y$，那么 $g(y)$ 就完全没有定义。换句话说，为了让 $g$ 能定义出来，$f$ 必须既是单射又是满射。

我们不妨慢慢来。我们把问题分成两部分：给定函数 $f\colon A \to B$，何时存在一个函数 $g\colon B \to A$，使得 $g(f(x)) = x$？这样的 $g$ “撤销”了 $f$ 的作用，称为 $f$ 的一个\emph{左逆}。其次，何时存在一个函数 $h\colon B \to A$，使得 $f(h(y)) = y$？这样的 $h$ 称为 $f$ 的一个\emph{右逆}——$f$ “撤销”了 $h$ 的作用。
\end{explain}

\begin{prop}
若 $f\colon A \to B$ 是单射且 $A \neq \emptyset$，则存在 $f$ 的一个\emph{左逆} $g\colon B \to A$，使得对所有 $x \in A$ 都有 $g(f(x)) = x$。
\end{prop}

\begin{proof}
假设 $f\colon A \to B$ 是单射且 $A \neq \emptyset$。考虑 $y \in B$。
若 $y \in \ran{f}$，则存在 $x \in A$ 使得 $f(x) = y$。因为 $f$ 是单射，这样的 $x \in A$ 是唯一的。于是我们可以定义：$g(y) = x$，即 $g(y)$ 是使得 $f(x) = y$ 的“那个”$x \in A$。
若 $y \notin \ran{f}$，我们可以将其映射到任意一个 $a \in A$。因此，我们可以选取一个 $a \in A$，并如下定义 $g\colon B \to A$：
\[
g(y) = \begin{cases}
    x & \text{若 $f(x) = y$}\\
    a & \text{若 $y \notin \ran{f}$。}
\end{cases}
\]
它对所有 $y \in B$ 都有定义，因为对于每个 $y \in \ran{f}$，都恰好存在一个 $x \in A$ 使得 $f(x) = y$。根据定义，如果 $y = f(x)$，那么 $g(y) = x$，即 $g(f(x)) = x$。
\end{proof}

\begin{prob}
证明：若 $f\colon A \to B$ 有一个左逆 $g$，则 $f$ 是单射。
\end{prob}

\begin{prop}
    若 $f\colon A \to B$ 是满射，则存在 $f$ 的一个\emph{右逆} $h\colon B \to A$，使得对所有 $y \in B$ 都有 $f(h(y)) = y$。
\end{prop}

\begin{proof}
假设 $f\colon A \to B$ 是满射。考虑 $y \in B$。由于 $f$ 是满射，存在 $x_y \in A$ 使得 $f(x_y) = y$。于是我们可以定义：$h(y) = x_y$，即对每个 $y \in B$，我们选择某个满足 $f(x) = y$ 的 $x \in A$；由于 $f$ 是满射，总是至少有一个可供选择。\footnote{由于 $f$ 是满射，对每个 $y \in B$，集合 $\Setabs{x}{f(x) = y}$ 非空。我们对 $h$ 的定义要求从这些集合中的每一个里选出一个 $x$。这种选择总是可行的其实并不显然——做出这些选择的可能性只是被直接作为公理假定的。换句话说，这个命题假定了所谓的选择公理，这个问题我们将在 \oliflabeldef{sth:choice::chap}{\olref[sth][choice][]{chap} 中再讨论}{略过不提}。然而，在许多具体情况下，例如当 $A = \Nat$ 或为有限集，或者当 $f$ 是双射时，并不需要选择公理。（特别地，当 $f$ 是双射时，对每个 $y \in B$，集合 $\Setabs{x}{f(x) = y}$ 恰好有一个元素，因此无需做出选择。）} 根据定义，如果 $x = h(y)$，那么 $f(x) = y$，即对任意 $y \in B$，$f(h(y)) = y$。
\end{proof}

\begin{prob}
证明：若 $f\colon A \to B$ 有一个右逆~$h$，则 $f$ 是满射。
\end{prob}

\begin{explain}
  结合前面证明中的思想，我们现在知道每个双射都有逆，也就是说，存在一个函数，它同时是~$f$ 的左逆和右逆。
\end{explain}

\begin{prop}\ollabel{prop:bijection-inverse}
若 $f\colon A \to B$ 是双射，则存在一个函数~$f^{-1}\colon B \to A$，使得对所有 $x \in A$，$f^{-1}(f(x)) = x$且对所有 $y \in B$，$f(f^{-1}(y)) = y$。
\end{prop}

\begin{proof}
练习。
\end{proof}

\begin{prob}
证明 \olref[sfr][fun][inv]{prop:bijection-inverse}。你必须定义~$f^{-1}$，证明它是一个函数，并证明它是~$f$ 的逆，即对所有 $x \in A$ 和 $y \in B$，$f^{-1}(f(x)) = x$和$f(f^{-1}(y)) = y$。
\end{prob}

\begin{explain}
还有一种稍具一般性的提取逆的方式。我们在 \olref[kin]{sec} 中看到，每个函数 $f$ 都通过对所有 $x \in A$ 令 $f'(x) = f(x)$ 诱导出一个满射 $f' \colon A \to \ran{f}$。显然，如果 $f$ 是单射，那么 $f'$ 是双射，从而根据 \olref{prop:bijection-inverse} 它有唯一的逆。出于一种轻微的记法滥用，我们有时直接把 $f'$ 的逆称为“$f$ 的逆”。
\end{explain}

\begin{prop}\ollabel{prop:left-right}%
  证明：若 $f\colon A \to B$ 有一个左逆~$g$ 和一个右逆~$h$，则 $h = g$。
\end{prop}

\begin{proof}
  练习。
\end{proof}

\begin{prob}
  证明 \olref[sfr][fun][inv]{prop:left-right}。
\end{prob}

\begin{prop}\ollabel{prop:inverse-unique}
每个函数~$f$ 至多有一个逆。
\end{prop}

\begin{proof}
  设 $g$ 和 $h$ 都是~$f$ 的逆。那么特别地，$g$~是~$f$ 的一个左逆，而 $h$~是一个右逆。根据 \olref{prop:left-right}，$g = h$。
\end{proof}

\end{document}